\documentclass[10pt,a4paper]{article}
\usepackage[hmargin=1.85cm,vmargin=1.55cm]{geometry}
\usepackage{fontspec}
\usepackage{booktabs,tabularx}
\usepackage{titlesec}
\usepackage{enumitem}
\setlength{\parindent}{0pt}
\setlength{\parskip}{3pt}
\setlength{\emergencystretch}{1em}
\titlespacing*{\section}{0pt}{7pt}{2pt}
\titlespacing*{\subsection}{0pt}{5pt}{2pt}
\setlist[enumerate]{nosep, topsep=3pt}

\begin{document}
\begin{center}
{\LARGE\bfseries Assignment 006: Vector\par}
\medskip
Zheng Yifeng \textbar{} Student ID: 3250105856
\end{center}

\section{Introduction}
This assignment implements a dynamic array as a C++ class template.
It was built with g++ 13.3.0 (C++17) on Ubuntu 24.04.4 LTS under WSL2.

\section{Program Design}

\subsection{Data Layout}
\texttt{Vector} uses a contiguous array to store data,which increase its capacity like a dynamic table.
\texttt{m\_pElements} points to the array, \texttt{m\_nSize} counts the elements in use, and \texttt{m\_nCapacity} counts the allocated slots.

When an append needs more space, \texttt{inflate()} doubles the capacity, copies the elements into a new array, frees the old array, and updates the pointer.
If the capacity is zero, it allocates one slot.

\subsection{Element Access}
Both \texttt{operator[]} and \texttt{at()} return an element by reference.
\texttt{operator[]} uses the index directly.
\texttt{at()} throws \texttt{std::out\_of\_range} if the index is outside \texttt{[0, m\_nSize)}.

\subsection{Appending an Element}
\texttt{push\_back()} writes at the current size and increases the size by one.
If the array is full, it saves the value before growing the array.
This keeps \texttt{v.push\_back(v[0])} safe when growth releases the array that \texttt{v[0]} refers to,that avoid memory leak.

\section{Testing}
Run from the assignment directory. All test inputs are in \texttt{main.cpp}.
\begin{verbatim}
g++ -std=c++17 -Wall -Wextra -Wpedantic code/main.cpp -o ../vector06-test
../vector06-test
\end{verbatim}

\small
\renewcommand{\arraystretch}{1.13}
\begin{tabularx}{\linewidth}{@{}>{\raggedright\arraybackslash}p{0.23\linewidth}*{2}{>{\raggedright\arraybackslash}X}@{}}
\toprule
\textbf{Test} & \textbf{Input or operation} & \textbf{Observed result} \\
\midrule
Default construction and growth & Start empty and append 100 integers from 0 to 990, in steps of 10 & Size is 100 and every value is preserved \\
Access and bounds & Write through both access functions and call \texttt{at(-1)} and \texttt{at(size())} & Writes are visible and both invalid indices throw \texttt{std::out\_of\_range} \\
Copy and destruction & Copy two strings, change the copy, and let both objects go out of scope & Original first string remains \texttt{hello} \\
Clear and reuse & Clear twice, then append 7 & Empty after clearing. Size is 1 and the first value is 7 after appending \\
Zero and negative sizes & Construct with 0 and append 8. Construct with -1 & Zero-size vector accepts 8. Negative size throws \texttt{std::invalid\_argument} \\
Positive size & Construct with 3 and check initial values. Write two slots & Three initial zeros. First and last values become 11 and 33 \\
Empty copy & Copy an empty vector and append 21 to the copy & Original stays empty. Copy holds 21 \\
Append an existing element & Call \texttt{v.push\_back(v[0])} on a full vector containing \texttt{hello} & The two elements are both \texttt{hello} \\
\bottomrule
\end{tabularx}
\normalsize

\section{Conclusion}
All eight test groups passed. The Vector class is built correctly.

\section{AI Use in the Lab}
\begin{enumerate}
  \item debug the program and fix small errors.
  \item fix memory leaks and learn how to manage dynamically allocated memory.
  \item find \texttt{v.push\_back(v[0])} special case, and handle it.
  \item generate test cases and part of the test program.
\end{enumerate}
\end{document}
